\chapter{Deterministic Gaussian Projection and Point-Rule Foundations}
\label{ch:bf-highdim-gaussian-point-rule-foundations}
\providecommand{\widehatell}{\widehat{\ell}}
\providecommand{\dotwidehatell}{\dot{\widehat{\ell}}}
\providecommand{\dd}{\,\mathrm d}

\section{Introduction And Scope}
\label{sec:bf-hd-gq-intro}

High-dimensional nonlinear filtering forces a choice of computational object.
The exact Bayesian filter propagates a full conditional law through time, but in
nonlinear models that law usually does not remain available in a finite closed
form.  In low dimension one may still tolerate heavy deterministic quadrature or
a richer non-Gaussian representation.  In high dimension, however, both cost and
structural complexity escalate quickly.  This chapter studies one narrower lane:
a deterministic Gaussian-surrogate filter that carries a single Gaussian object,
uses explicit moment engines to approximate the nonlinear expectations it needs,
and prepares the rule-family language that the next chapter will specialize to
fixed sparse-grid quadrature.

This chapter imports the exact-target and same-scalar discipline from
Chapter~\ref{ch:bf-highdim-nonlinear-foundations}.  Its task is not yet to build
the full sparse-grid algorithm.  It first fixes the deterministic Gaussian
carried object, the Gaussian projection contract, and the point-rule family
comparisons needed before sparse-grid specialization.  The next chapter,
Chapter~\ref{ch:bf-highdim-sparse-grid-fixed-cloud}, then turns this shared
Gaussian-rule foundation into the explicit fixed-cloud sparse-grid lane.

The running quadratic-observation cell from
Chapter~\ref{ch:bf-highdim-nonlinear-foundations} reappears here as a compact
shared object, but the present chapter stays at the level of Gaussian carried
objects and deterministic moment engines.  Low-dimensional cloud construction,
value recursion, and the fixed-cloud scalar belong to the sparse-grid chapter
that follows.

The central limitation should be fixed at the outset.  A deterministic Gaussian
lane does not claim exact nonlinear filtering, exact nonlinear likelihood
evaluation, or faithful preservation of arbitrary multimodality, skewness, or
high-order global interaction structure.  Its purpose is narrower: to ask
whether one carried Gaussian, together with a declared deterministic moment
engine, can provide a coherent and scientifically useful approximate likelihood
lane in the regime under study.

\paragraph{Background assumed.}
The reader is assumed to know multivariate Gaussian notation, state-space-model
notation, Jacobian-style derivative notation, and the use of a Cholesky factor
for a positive-definite covariance matrix.  Sparse-grid merging, same-scalar
branch identities, and fixed-cloud derivative obligations are defined where they
first become necessary.

\subsection{Problem Setting And Deterministic Gaussian Lanes}
\label{sec:bf-hd-gq-problem-lane}

The scientific problem is high-dimensional nonlinear filtering.  At time
\(t\), the exact Bayesian filter propagates and updates the full conditional law
\(p_\theta(x_t\mid y_{1:t})\).  In nonlinear models that law is generally not
available in a finite-dimensional closed form, and in high dimension it becomes
unrealistic to treat exact function-valued propagation as the practical
computational object.  The deterministic Gaussian family therefore does not
attempt to rescue exactness by notation.  It selects one approximate lane whose
strengths and boundaries can both be stated explicitly.

Within this family, the carried object is a single Gaussian surrogate
\(\calN(m_t,P_t)\).  Different deterministic point-rule lanes then differ not in
what object is carried, but in how the nonlinear Gaussian moments needed by the
projection are approximated.  This chapter fixes that shared carried-object and
moment-contract language before the next chapter specializes it to the fixed
sparse-grid lane.

\subsection{The Deterministic Gaussian Lane In One Page}
\label{sec:bf-hd-gq-one-page}

The deterministic Gaussian family studied here is defined by two deliberate
approximations and one downstream derivative constraint.
\begin{enumerate}[leftmargin=0.28in]
  \item \textbf{Gaussian carried state.}  Instead of carrying the full filtering
  law, the method carries one Gaussian summary \(\calN(m_t,P_t)\).
  \item \textbf{Deterministic moment engine.}  The nonlinear Gaussian moments that
  define the prediction and observation projection are approximated by a fixed
  point-rule family in standardized coordinates.
  \item \textbf{Same-scalar derivative discipline.}  Any later reported gradient must
  be the derivative of one declared approximate scalar built from that same
  deterministic moment engine.
\end{enumerate}

The exact object, the carried surrogate, and the eventual approximate scalar are
therefore different in kind:
\begin{equation}
\begin{aligned}
  \text{exact object:}
  &\quad p_\theta(x_t\mid y_{1:t}), \\
  \text{carried surrogate:}
  &\quad \calN(m_t,P_t), \\
  \text{declared approximate scalar:}
  &\quad \widehat\ell_T^{(M)}(\theta)
  \text{ for a chosen deterministic moment engine } M.
\end{aligned}
  \label{eq:bf-hd-gq-exact-vs-approximate-summary}
\end{equation}
This distinction is the governing contract of the deterministic Gaussian block.
The next chapter will instantiate one specific engine, fixed sparse-grid
quadrature, within this larger family.

\subsection{Approximation Hierarchy And Reading Guide}
\label{sec:bf-hd-gq-hierarchy}

The deterministic Gaussian side of the expanded block unfolds in two chapter
movements.  First, the exact filtering recursion and the Gaussian projection that
replaces it are fixed formally.  Second, nearby deterministic point-rule families
are compared at the level of their moment engines.  Only after that shared
foundation is explicit does the next chapter specialize to sparse-grid
construction, merged clouds, the filtering value path, and the fixed-cloud
scalar.

The main coordinates are:
\begin{align}
  \xi &\in \R^{n_x},
  &\xi&\sim\calN(0,I_{n_x})
  &&\hbox{standard quadrature coordinate},
  \label{eq:bf-hd-gq-standard-coordinate}\\
  x &= m+C\xi,
  &CC^\top&=P
  &&\hbox{physical state coordinate under a carried Gaussian},
  \label{eq:bf-hd-gq-physical-coordinate}\\
  a &= f_\theta(x),
  &&
  &&\hbox{transition image used for prediction},
  \label{eq:bf-hd-gq-transition-coordinate}\\
  z &= h_\theta(x),
  &&
  &&\hbox{observation image used for the innovation}.
  \label{eq:bf-hd-gq-observation-coordinate}
\end{align}
A deterministic point-rule lane builds a finite rule in \(\xi\)-space and then
places that rule into physical state space through the current Gaussian factor.
That sentence is the bridge from the present foundations chapter to the
sparse-grid chapter that follows.

\section{Exact Filtering And Gaussian Projection}
\label{sec:bf-hd-gq-exact-filtering}

Use the additive-noise state-space model
\begin{align}
  X_t &= f_\theta(X_{t-1})+\eta_t,
  &
  \eta_t&\sim\calN(0,Q_\theta),
  \label{eq:bf-hd-gq-state-model}\\
  Y_t &= h_\theta(X_t)+\varepsilon_t,
  &
  \varepsilon_t&\sim\calN(0,R_\theta),
  \label{eq:bf-hd-gq-observation-model}
\end{align}
with \(X_t\in\R^{n_x}\), \(Y_t\in\R^{n_y}\), and parameter
\(\theta\in\R^p\).  The exact Bayesian prediction and update are
\begin{align}
  p_\theta(x_t\mid y_{1:t-1})
  &=
  \int p_\theta(x_t\mid x_{t-1})
       p_\theta(x_{t-1}\mid y_{1:t-1})\dd x_{t-1},
  \label{eq:bf-hd-gq-exact-prediction}\\
  p_\theta(x_t\mid y_{1:t})
  &=
  \frac{
  g_\theta(y_t\mid x_t)
  p_\theta(x_t\mid y_{1:t-1})}
  {\int g_\theta(y_t\mid x_t)
  p_\theta(x_t\mid y_{1:t-1})\dd x_t}.
  \label{eq:bf-hd-gq-exact-update}
\end{align}
In nonlinear models, these integrals usually do not close in finite-dimensional
form.  Deterministic Gaussian filters replace those exact densities by a Gaussian
moment projection.

\subsection{Object Map And Deterministic Gaussian Conventions}
\label{subsec:bf-hd-gq-object-map}

The deterministic Gaussian family uses the following common objects.
\begin{center}
\small
\begin{tabular}{p{0.20\linewidth}p{0.44\linewidth}p{0.22\linewidth}}
\toprule
symbol & meaning & shape \\
\midrule
\(m_t,P_t\) & carried Gaussian mean and covariance after update & \(n_x\), \(n_x\times n_x\) \\
\(m_t^-,P_t^-\) & predictive Gaussian mean and covariance before update & \(n_x\), \(n_x\times n_x\) \\
\(C_t,C_t^-\) & covariance factors with \(P_t=C_tC_t^\top\), \(P_t^-=C_t^-(C_t^-)^\top\) & \(n_x\times n_x\) \\
\(\chi_t^{(r)}\) & deterministic physical rule point & \(n_x\) \\
\(z_t^{(r)}\) & observation image of a rule point & \(n_y\) \\
\(\bar z_t,S_t,C_{xz,t}\) & predicted observation mean, innovation covariance, cross-covariance & \(n_y\), \(n_y\times n_y\), \(n_x\times n_y\) \\
\(v_t,K_t\) & innovation residual and Gaussian projection gain & \(n_y\), \(n_x\times n_y\) \\
\bottomrule
\end{tabular}
\end{center}

Three implementation conventions are fixed chapter-wide.
\begin{enumerate}[leftmargin=0.28in]
  \item \textbf{Factor convention.}  Every covariance factor is the lower
  triangular Cholesky factor with positive diagonal, written \(P=CC^\top\),
  whenever the declared branch is valid.
  \item \textbf{Solve convention.}  Formulas may use \(S_t^{-1}\) for compact
  mathematics, but the implementation contract is to compute such actions by
  linear solves against the declared factor of \(S_t\), not by forming an
  explicit inverse.
  \item \textbf{Noise convention.}  Additive process noise enters analytically
  through \(Q_\theta\) in predictive moments and additive observation noise
  through \(R_\theta\) in innovation moments.  This chapter does not use state
  augmentation for these additive noises.
\end{enumerate}

\subsubsection{Gaussian Projection From Moments}
\label{sec:bf-hd-gq-projection}

Assume the predictive state is represented by
\begin{equation}
  X_t^-\sim\calN(m_t^-,P_t^-).
  \label{eq:bf-hd-gq-predictive-gaussian}
\end{equation}
Let
\begin{equation}
  Z_t=h_\theta(X_t^-)
  \label{eq:bf-hd-gq-predicted-observation}
\end{equation}
be the noiseless predicted observation.  The Gaussian projection needs
\begin{align}
  \bar z_t &= \E[Z_t],
  \label{eq:bf-hd-gq-zbar}\\
  S_t &= R_\theta+\E[(Z_t-\bar z_t)(Z_t-\bar z_t)^\top],
  \label{eq:bf-hd-gq-S}\\
  C_{xz,t}&=\E[(X_t^- -m_t^-)(Z_t-\bar z_t)^\top].
  \label{eq:bf-hd-gq-Cxz}
\end{align}
Given these moments, define
\begin{align}
  v_t&=y_t-\bar z_t,
  &
  K_t&=C_{xz,t}S_t^{-1},
  \label{eq:bf-hd-gq-v-K}\\
  m_t&=m_t^-+K_tv_t,
  &
  P_t&=P_t^- - K_tS_tK_t^\top.
  \label{eq:bf-hd-gq-projection-update}
\end{align}
The approximation has two layers.  First, these moments may be approximated by a
deterministic point rule.  Second, the resulting posterior information is
compressed back to one Gaussian.

\subsubsection{How This Matches The Jia--Xin--Cheng Gaussian Approximation Block}
\label{subsec:bf-hd-gq-jia-block}

Jia, Xin, and Cheng first write the Gaussian approximation filter before they
introduce sparse grids.  In their additive-noise setting, the sparse-grid part of
the method answers only how to approximate the Gaussian expectations needed for
\(m_t^-\), \(P_t^-\), \(\bar z_t\), \(S_t\), and \(C_{xz,t}\).  It does not
change the fact that the filter stores one Gaussian pair \((m_t,P_t)\).

In BayesFilter notation, that bridge has the generic form
\begin{equation}
  \E[F(\xi)]
  \quad\leadsto\quad
  \mathcal Q_M[F]
  \quad\leadsto\quad
  \sum_{r} w_r^{(M)} F(\xi_r^{(M)}),
  \label{eq:bf-hd-gq-expectation-to-rule}
\end{equation}
where \(M\) denotes the chosen deterministic point-rule family.  The next
chapter specializes this generic expectation-to-rule replacement to the fixed
sparse-grid cloud and its stored scalar contract.

\section{Deterministic Point-Rule Families}
\label{sec:bf-hd-gq-rule-families}

A deterministic Gaussian-surrogate lane is defined not only by the carried
Gaussian object but also by the moment engine it uses.  The key families for the
expanded high-dimensional block are:
\begin{itemize}[leftmargin=0.28in]
  \item local linearization (EKF-style moment closure),
  \item low-order sigma-point rules (UKF/CKF-style closures),
  \item dense Gauss--Hermite / tensor-product rules,
  \item sparse-grid rules that the next chapter will specialize.
\end{itemize}
The decisive question is therefore not whether all these methods are “Gaussian,”
but what deterministic moment engine they use and what that engine implies for
the approximate scalar that later chapters may differentiate.

The standard sigma-point interface already exists elsewhere in the monograph.
Chapter~\ref{ch:bf-sigma-point-filters} gives the generic sigma-point
moment-matching interface, while Chapter~\ref{ch:bf-square-root-sigma-point}
clarifies square-root backend contracts.  The purpose here is not to reteach all
sigma-point filters, but to position the deterministic point-rule family as the
shared foundations of the sparse-grid specialization.

\paragraph{Rule-family comparison discipline.}
For a deterministic point rule with nodes \(\xi_r^{(M)}\), mean weights
\(w_{r,m}^{(M)}\), and covariance/cross-covariance weights \(w_{r,c}^{(M)}\),
the family-level operator has the form
\begin{equation}
  \bar z_t^{(M)}
  \approx
  \sum_r w_{r,m}^{(M)} h_\theta(m_t^-+C_t^-\xi_r^{(M)}),
  \label{eq:bf-hd-gq-family-zbar}
\end{equation}
\begin{equation}
  S_t^{(M)}
  \approx
  R_\theta+
  \sum_r w_{r,c}^{(M)}\Delta_{t,r}^{(M)}\Delta_{t,r}^{(M)\top},
  \label{eq:bf-hd-gq-family-S}
\end{equation}
with the matching cross-covariance operator defined analogously.  The one-weight
specialization covers symmetric rules such as CKF-like or fixed sparse-grid
clouds.  Standard scaled-UKF comparisons generally require distinct mean and
covariance weights, so rule-family comparisons must preserve that distinction.

\section{Exports To Sparse-Grid Specialization And Later Chapters}
\label{sec:bf-hd-gq-exports}

This chapter exports four things forward.
\begin{enumerate}[label=(\arabic*), leftmargin=0.28in]
  \item the exact-target versus Gaussian-carried-object distinction,
  \item the Gaussian projection contract and its innovation objects,
  \item the deterministic point-rule family language needed before sparse-grid
  specialization,
  \item the rule-family comparison discipline needed later when sparse-grid,
  low-order sigma-point, and dense GHQ lanes are compared.
\end{enumerate}

The next chapter, Chapter~\ref{ch:bf-highdim-sparse-grid-fixed-cloud}, begins
where this one stops: it instantiates one specific deterministic point-rule lane,
fixed sparse-grid quadrature, and develops the low-dimensional cloud geometry,
Smolyak structure, filtering value path, and fixed-cloud scalar that this shared
foundations chapter has intentionally deferred.
